\documentclass[11pt,a4paper]{article}
\usepackage[T1]{fontenc}
\usepackage[utf8]{inputenc}
\usepackage{amsmath,amssymb,amsthm}
\usepackage[margin=1in]{geometry}
\usepackage[colorlinks=true,linkcolor=blue,urlcolor=blue]{hyperref}
\theoremstyle{plain}
\newtheorem{theorem}{Theorem}
\newtheorem{lemma}[theorem]{Lemma}
\newtheorem{proposition}[theorem]{Proposition}
\newtheorem{corollary}[theorem]{Corollary}
\theoremstyle{definition}
\newtheorem{remark}[theorem]{Remark}

\title{The empirical recurrence for OEIS A234439, proved by transfer matrix}
\author{Adrian Perez Fontelles\\ \small Independent researcher}
\date{31 August 2026}
\begin{document}
\maketitle

\begin{abstract}
OEIS A234439 counts $5$-column arrays over $\{0,\dots,2\}$ in which every
$2\times2$ subblock has its diagonal sum differing from its antidiagonal sum by
exactly $1$, and it
carries an empirical recurrence of order $45$ contributed by R.~H.~Hardin. It is true,
and it is not an empirical matter: the condition constrains only consecutive pairs of rows,
so the rows are the vertices of a finite digraph and an $(n+1)$-row array is a walk of
length $n$. Hence $a(n)=\mathbf{1}^{\!\top}M^{n}\mathbf{1}$ on $S=243$ vertices, which
is $C$-finite. Writing $q$ for the characteristic polynomial of the conjectured recurrence,
the residual is $\mathbf{1}^{\!\top}M^{\,n-45}q(M)\mathbf{1}$, so the recurrence
holds beyond a threshold exactly when that vanishes from then on --- and since the sequence
of those values obeys the monic recurrence given by the characteristic polynomial of $M$,
$S$ consecutive zeros settle every later one. The whole test is finite and exact.
\end{abstract}

\noindent\small 2020 Mathematics Subject Classification. 05A15, 11B37, 15B36.\normalsize

\section{The sequence and the conjecture}

OEIS A234439 is ``Number of (n+1) X (4+1) 0..2 arrays with every 2 X 2 subblock having its diagonal sum differing from its antidiagonal sum by 1 (constant-stress 1 X 1 tilings).''. It has offset
$1$ and begins
\[
1680, 13676, 108404, 896112, 7165920, 59991972, 483370700,\ \dots
\]
The entry carries the following comment:
\begin{quote}\small
Empirical: a(n) = 19*a(n-1) +174*a(n-2) -5225*a(n-3) -579*a(n-4) +557372*a(n-5) -1736348*a(n-6) -28633611*a(n-7) +147097160*a(n-8) +701879450*a(n-9) -5215474532*a(n-10) -6448644414*a(n-11) +90991389225*a(n-12) -13155453655*a(n-13) -911859068299*a(n-14) +785985063499*a(n-15) +5750480317192*a(n-16) -7719561708652*a(n-17) -23967944021653*a(n-18) +42175716887542*a(n-19) +67283824235584*a(n-20) -150636045987035*a(n-21) -125392746919991*a(n-22) +373892366562458*a(n-23) +141881972068085*a(n-24) -664306804785765*a(n-25) -55299787514109*a(n-26) +857132068732641*a(n-27) -106700542512030*a(n-28) -806720994198218*a(n-29) +225446251980190*a(n-30) +551556767133038*a(n-31) -222526751973132*a(n-32) -270104055145308*a(n-33) +138509404758088*a(n-34) +92036096585608*a(n-35) -57527289381280*a(n-36) -20591630958048*a(n-37) +15900557729248*a(n-38) +2635188748768*a(n-39) -2803132131904*a(n-40) -102351217472*a(n-41) +284333800832*a(n-42) -15435123840*a(n-43) -12584140800*a(n-44) +1510096896*a(n-45).
\end{quote}
As of the ``Last modified'' line on the live entry (Jun 20 2022 00:33:10, revision 6) the statement is
still recorded as empirical, and nothing on the entry records it as settled either way.

Written out, the claim is
\begin{equation}\label{eq:conj}
a(n)\;=\;19\,a(n-1) + 174\,a(n-2) - 5225\,a(n-3) - 579\,a(n-4) + 557372\,a(n-5) - 1736348\,a(n-6) - 28633611\,a(n-7) + 147097160\,a(n-8) + 701879450\,a(n-9) - 5215474532\,a(n-10) - 6448644414\,a(n-11) + 90991389225\,a(n-12) - 13155453655\,a(n-13) - 911859068299\,a(n-14) + 785985063499\,a(n-15) + 5750480317192\,a(n-16) - 7719561708652\,a(n-17) - 23967944021653\,a(n-18) + 42175716887542\,a(n-19) + 67283824235584\,a(n-20) - 150636045987035\,a(n-21) - 125392746919991\,a(n-22) + 373892366562458\,a(n-23) + 141881972068085\,a(n-24) - 664306804785765\,a(n-25) - 55299787514109\,a(n-26) + 857132068732641\,a(n-27) - 106700542512030\,a(n-28) - 806720994198218\,a(n-29) + 225446251980190\,a(n-30) + 551556767133038\,a(n-31) - 222526751973132\,a(n-32) - 270104055145308\,a(n-33) + 138509404758088\,a(n-34) + 92036096585608\,a(n-35) - 57527289381280\,a(n-36) - 20591630958048\,a(n-37) + 15900557729248\,a(n-38) + 2635188748768\,a(n-39) - 2803132131904\,a(n-40) - 102351217472\,a(n-41) + 284333800832\,a(n-42) - 15435123840\,a(n-43) - 12584140800\,a(n-44) + 1510096896\,a(n-45)
\end{equation}
for all $n>45$, which is every $n$ at which a recurrence of order $45$ can be
stated.

\section{The rows form a digraph}

In a $2\times2$ subblock with top row $(r_j,r_{j+1})$ and bottom row $(s_j,s_{j+1})$ the
diagonal sum is $r_j+s_{j+1}$ and the antidiagonal sum is $r_{j+1}+s_j$, so the entry's
condition on that subblock reads
\[
\bigl|\,r_j+s_{j+1}-r_{j+1}-s_j\,\bigr|\;=\;1.
\]
Every such condition involves two consecutive rows only.

Let $V$ consist of all of $\{0,\dots,2\}^{5}$, and put $S=|V|=243$. For $r,s\in V$ declare $r\to s$ an edge
when $\bigl|r_j+s_{j+1}-r_{j+1}-s_j\bigr|=1$ for every $j$ with
$1\le j\le5-1$. Let $M\in\{0,1\}^{S\times S}$ be the adjacency matrix and $\mathbf{1}$
the all-ones vector.

\begin{lemma}\label{lem:walk}
The arrays counted by the entry with $n+1$ rows are exactly the walks
$r^{(0)}\to\cdots\to r^{(n)}$ of length $n$ in this digraph, so
$a(n)=\mathbf{1}^{\!\top}M^{n}\mathbf{1}$.
\end{lemma}

\begin{proof}
An array with $n+1$ rows and $5$ columns is a sequence $r^{(0)},\dots,r^{(n)}$ of
rows. Each of its $2\times2$ subblocks lies in two vertically adjacent rows
$r^{(i)},r^{(i+1)}$ and two horizontally adjacent columns $j,j+1$, so the array satisfies
the entry's condition on all subblocks precisely when every consecutive pair of rows is an
edge; and the row-internal condition, where present, is exactly membership of $V$. Summing
the walk counts $\bigl(M^{n}\bigr)_{r,s}$ over all $r$ and $s$ gives
$\mathbf{1}^{\!\top}M^{n}\mathbf{1}$.
\end{proof}

\section{The criterion}

Let $q(t)=t^{45}-\sum_i c_i t^{45-i}$ be the characteristic polynomial of
\eqref{eq:conj}.

\begin{lemma}\label{lem:crit}
For $n\ge45$,
$a(n)-\sum_i c_i a(n-i)=\mathbf{1}^{\!\top}M^{\,n-45}q(M)\mathbf{1}$. Writing
$u_j=\mathbf{1}^{\!\top}M^{j}q(M)\mathbf{1}$, the recurrence therefore holds for every
$n>t$ exactly when $u_j=0$ for every $j>t-45$. Moreover $(u_j)_{j\ge0}$ satisfies the
monic linear recurrence whose characteristic polynomial is that of $M$, of degree $S$; so
if $u_j$ vanishes for $S$ consecutive indices it vanishes for every larger index, and the
last nonzero $u_j$ determines $t$ exactly.
\end{lemma}

\begin{proof}
By Lemma~\ref{lem:walk}, $a(n-i)=\mathbf{1}^{\!\top}M^{n-i}\mathbf{1}$, so
$a(n)-\sum_i c_i a(n-i)
=\mathbf{1}^{\!\top}M^{\,n-45}\bigl(M^{45}-\sum_i c_i M^{45-i}\bigr)
\mathbf{1}$, which is the identity. Putting $j=n-45$ gives the equivalence. For the
last claim, let $\chi$ be the characteristic polynomial of $M$; by Cayley--Hamilton
$\chi(M)=0$, so applying $\chi$ to the shift operator annihilates $(u_j)$, and $\chi$ is
monic of degree $S$, which is exactly the statement that $S$ consecutive zeros propagate
forward.
\end{proof}

\section{The computation}

\begin{theorem}
The recurrence \eqref{eq:conj} holds for every $n>45$. This is the statement of
Section~1.
\end{theorem}

\begin{proof}
The vector $q(M)\mathbf{1}\in\mathbb{Z}^{S}$ was formed by $45$ matrix--vector
products in exact integer arithmetic and the values $u_j$ evaluated for
$j=0,1,\dots$. They vanish from $j=45-45+1$ onwards, and the run was continued past
$S$ consecutive zeros, which by Lemma~\ref{lem:crit} settles every larger $j$. Hence the
recurrence holds for every $n>45$.
\end{proof}

\section{Verification}

Three checks, each independent of the algebra above.

First, the digraph was built directly from the entry's stated condition and the walk counts
$\mathbf{1}^{\!\top}M^{n}\mathbf{1}$ compared against the entry's DATA at every one of
the $16$ published indices; the values agree exactly. That is what ties the digraph to
the sequence: a model reproducing only some published terms would be the wrong model and
was rejected wherever it occurred.

Second, the recurrence \eqref{eq:conj} was evaluated on those published terms in exact
integer arithmetic, with no matrices involved, and vanishes wherever it applies.

Third, the vanishing test of Lemma~\ref{lem:crit} was carried past $S=243$ consecutive
zeros rather than to a sampled prefix, in exact integer arithmetic.

\begin{thebibliography}{9}
\bibitem{oeis} The OEIS Foundation, \emph{The On-Line Encyclopedia of Integer
Sequences}, \url{https://oeis.org}, 2026. Sequence A234439.
\bibitem{stanley} R.~P.~Stanley, \emph{Enumerative Combinatorics}, Volume 1, 2nd ed.,
Cambridge University Press, 2011. (Transfer-matrix method, Section 4.7.)
\bibitem{fs} P.~Flajolet and R.~Sedgewick, \emph{Analytic Combinatorics}, Cambridge
University Press, 2009.
\bibitem{horn} R.~A.~Horn and C.~R.~Johnson, \emph{Matrix Analysis}, 2nd ed., Cambridge
University Press, 2013.
\end{thebibliography}

\end{document}
